; DTXManiaAI 第5段の確認用譜面（拍線 / 0xC1 / 0xC2 / フィルイン / サビ / 歓声 / SE）
;
; 使い方: 選曲画面で "Demo Features" を選び、演奏中に F1 で AUTO にすると全部自動で流れる。
; メトロノームを聞くときは Config で Metronome=ON、拍線を見るときは Dark=OFF にしておくこと。
;
;  小節 000-001  4/4。小節線 1 本＋拍線 3 本（拍線は小節線より暗い）
;  小節 002      小節長 0.75（3/4）。拍線 2 本
;  小節 003      小節長 2.00。拍線 7 本
;  小節 004      小節長を 1.00 に戻し、0xC1 拍線シフト。拍線が 1/2 小節ぶんずれて 2 本になる
;                （小節長は NX 同様、次に指定されるまで持続する）
;  小節 005-006  0xC2=02 で小節線・拍線を非表示（メトロノームは鳴り続ける）
;  小節 007      0xC2=01 で表示を戻す
;  小節 008-009  フィルイン区間（0x53=01→02）。叩くと青い星、区間最後のチップで波、終了時に歓声
;  小節 010-011  サビ区間（0x53=03→04）。チップファイアにボーナス花火が重なる
;  小節 012      SE チャンネル（0x61）。#WAV を置かない限り無音（パース確認用）
;  小節 013      余韻

#TITLE: Demo Features
#ARTIST: DTXManiaAI
#COMMENT: 第5段の確認用（拍線/フィルイン/サビ/歓声）。F1 で AUTO にすると全部見える。
#GENRE: Sample
#BPM: 120
#DLEVEL: 30

; --- 000: 4/4 の基本。拍線 3 本 ---
#00011: 01010101
#00013: 01000100

; --- 001: 同上 ---
#00111: 01010101
#00113: 01000100

; --- 002: 小節長 0.75（3/4）→ 拍線 2 本 ---
#00202: 0.75
#00211: 010101
#00213: 010001

; --- 003: 小節長 2.00 → 拍線 7 本 ---
#00302: 2.0
#00311: 0101010101010101
#00313: 0100010001000100

; --- 004: 0xC1 拍線シフト（小節の 1/2 位置にチップを置く）→ 拍線が右へずれて 2 本 ---
;      小節長は NX 同様「次の #xxx02 まで持続」するので、003 の 2.0 をここで 1.0 に戻す
#00402: 1.0
#004C1: 0001
#00411: 01010101
#00413: 01000100

; --- 005: 0xC2=02 で小節線・拍線を非表示 ---
#005C2: 02
#00511: 01010101
#00513: 01000100

; --- 006: 非表示のまま ---
#00611: 01010101
#00613: 01000100

; --- 007: 0xC2=01 で表示を戻す ---
#007C2: 01
#00711: 01010101
#00713: 01000100

; --- 008: フィルイン開始（0x53=01）。以降のヒットで青い星が出る ---
#00853: 01
#00811: 0101010101010101
#00812: 00010001
#00814: 01000000
#00815: 00000100

; --- 009: フィルイン終了（0x53=02、小節末）。最後のチップで波、終了時に歓声 ---
;      0x1F は歓声チップ。#WAV を定義していないのでスキンの Audience.ogg が鳴る
#0091F: 01
#00911: 0101010101010101
#00912: 01010101
#00917: 00000101
#00916: 00000001
#00953: 00000002

; --- 010: サビ開始（0x53=03）。チップファイアにボーナス花火が重なる ---
#01053: 03
#01011: 01010101
#01013: 01000100
#01016: 01000000

; --- 011: サビ終了（0x53=04） ---
#01111: 01010101
#01113: 01000100
#01119: 00000001
#01153: 00000004

; --- 012: SE チャンネル（パース確認用。#WAV が無いので無音） ---
#01261: 01
#01211: 01010101
#01213: 01000100

; --- 013: 余韻 ---
#01311: 01
#01313: 01
